\metatitle{Families of permutations}
\metadescription{Named families of permutations: pattern-avoidance classes, alternating permutations, derangements, and more.}


\section[namedPermutations]{Named families of permutations}

See the \hyperref[permutationPattern]{permutation patterns page}
for background on pattern containment and avoidance.

\begin{itemize}

\item \defin{Up-down}\label{permUpDown} permutations have, in one-line
notation, $\pi_1 \lt \pi_2 \gt \pi_3 \lt \pi_4 \gt \dotsb$.
\defin{Down-up} permutations have
$\pi_1 \gt \pi_2 \lt \pi_3 \gt \pi_4 \lt \dotsb$.
Each orientation is enumerated by the Euler numbers \oeis{A000111}.
If both orientations are included, then for $n\gt 1$ the count is twice an
Euler number, as in \oeis{A001250}.

\item \defin{Derangements} --- permutations with no fixed point, \oeis{A000166}.

\item \defin{Dominant} --- 132-avoiding.
Enumerated by the Catalan numbers, \oeis{A000108}.

\item \defin{Baxter} --- avoids the vincular patterns $2\text{-}41\text{-}3$
and $3\text{-}14\text{-}2$, see \cite{DulucqGuibert1998} and
\oeis{A001181}.

\item \defin{Boolean} --- 321-avoiding and 3412-avoiding, see
\cite{Tenner2007,GunawanPanRussellTenner2022x}.
A permutation is boolean if and only if the order ideal it defines in the
\hyperref[bruhatOrder]{Bruhat order} is a Boolean lattice,
see \cite[Thm.~4.3]{Tenner2007}.
Enumerated by odd-indexed Fibonacci numbers, \oeis{A001519}.

\item \defin{Fireworks} --- Same as 3-12 avoiding, see
\cite{ChouSetiabrata2025x}. This means that the initial terms of decreasing
runs are increasing.  Fireworks permutations are enumerated by the Bell
numbers \oeis{A000110}.

\item \defin{Fully commutative} --- 321-avoiding, see
\cite{BilleyJockuschStanley1993,GunawanPanRussellTenner2022x}.
Enumerated by the Catalan numbers, \oeis{A000108}.

\item A \defin{simsun permutation} is a permutation such that, after restricting
to entries $1,2,\dotsc,k$, it does not have a double descent, that is,
$\pi(i)\gt \pi(i+1) \gt \pi(i+2)$, see \cite{MaYeh2016}.
The number of simsun permutations in $\symS_n$ is the Euler number $E_{n+1}$;
see \oeis{A000111}.

\item \defin{King permutations} are permutations in which no two adjacent
entries differ by $1$ in one-line notation.
This can be described via mesh-pattern avoidance, see \oeis{A002464}.

\item \label{vexillary}
A permutation is \defin{vexillary} if it avoids the pattern 2143; see
\oeis{A005802}.

\item A permutation is \defin{Grassmann} (or \defin{Grassmannian}) if it has
at most one descent.
These are special cases of vexillary permutations.
There are $2^n-n$ Grassmannian permutations in $\symS_n$; see
\oeis{A000325}.  See
\hyperref[grassmannianPermutation]{the Schubert polynomial page} for more
details.

\item A permutation $\pi$ is \defin{bigrassmannian} if both $\pi$ and
$\pi^{-1}$ are Grassmannian.  The number of such permutations in $\symS_n$,
including the identity, is $1+\binom{n+1}{3}$; see \oeis{A050407}, with its
index shifted by two.

\item \defin{Separable} permutations are those which avoid $2413$ and $3142$,
see \oeis{A006318}.
For reference, see \cite{West1995}.

\item \defin{Skew-merged} permutations avoid $2143$ and $3412$. These are
exactly the ones formed as the union of an increasing sequence and a decreasing
sequence; see \cite[Thm.~2.9]{Stankova1994}, \oeis{A029759}.
See \cite{Atkinson1998} for exact enumeration and connection with RSK.

\item \defin{Flattened permutations}, where runs are in lex order.
The flattened permutations in $\symS_n$ are enumerated by the shifted Bell
number $B_{n-1}$; see \oeis{A000110} and
\cite[Prop.~16]{NabawandaRakotondrajaoBamunoba2020}.

\item \defin{Parity-alternating permutations} are permutations satisfying
$\pi(i)\equiv i\pmod 2$, see \cite{KebedeRakotondrajao2021}.
There are $\lfloor n/2\rfloor!\,\lceil n/2\rceil!$ such permutations in
$\symS_n$; see \oeis{A010551}.
This can be generalized to \defin{modulo-$k$ alternating permutations},
where we demand that $\pi(i) \equiv i \pmod k$.
For pattern avoidance in modulo-$k$ alternating permutations, see
\cite{AlexanderssonFufaGetachewQiu2022x}.


\item For a composition $(k_1,k_2,\dotsc,k_m)$ of $n$, the corresponding
\defin{Richardson permutation}, also called a \defin{layered permutation}, is
\[
  w_0^{(k_1)}\oplus w_0^{(k_2)}\oplus\dotsb\oplus w_0^{(k_m)},
\]
the \hyperref[directSumPermutations]{direct sum} of the longest permutations
$w_0^{(k_i)}\in\symS_{k_i}$; see \cite{MerzonSmirnov2015}.
Equivalently, layered permutations avoid both 231 and 312.  There are
$2^{n-1}$ layered permutations in $\symS_n$, one for every composition of
$n$; see \oeis{A000079}.

\item \defin{Wachs permutations} --- \oeis{A359039}.
These are the permutations in $\symS_n$ defined as
\[
 W_n \coloneqq
 \{\sigma\in\symS_n:|\sigma^{-1}(i)-\sigma^{-1}(i^*)|\leq 1
 \text{ for all }i\in[n-1]\},
\]
where $i^*$ is defined as
\[
 i^* \coloneqq \begin{cases}
 i-1 &\text{if $i$ is even}  \\
 i+1 &\text{if $i$ is odd and $i+1 \leq n$}  \\
 n & \text{otherwise.}
 \end{cases}
\]
These permutations are studied in \cite{BrentiSentinelli2022x}, and there are generalizations to
other types.

\end{itemize}
